\section{Deviations and irregularities of Study 3}
\label{app:deviations}
\label{app:dev-study3}

This appendix covers Study~3; the deviations of Studies~1 and~2 are in the Supplement (Sections~\ref{S-app:dev-study1} and~\ref{S-app:dev-study2}). The analysis plan of Study~3 (\texttt{ANALYSIS\_\allowbreak{}PLAN\_\allowbreak{}STUDY3.md}) was committed before any design ran on \texttt{fresh2/}; the data loader and the run driver refused to start unless it was committed unmodified. The analysis (\texttt{analysis/\allowbreak{}study3.py} with \texttt{analysis/\allowbreak{}study3\_\allowbreak{}spec.json}) was run once on the stored outputs. Times are UTC.

\subsection{The run}

\paragraph{How the invocations ended}
The run took \sthreeRunInvocations{} invocations of the driver, from \sthreeRunFirstDate{} \sthreeRunFirstStart{} to \sthreeRunLastDate{} \sthreeRunLastEnd. Of these, \sthreeRunNotComplete{} stopped early: \sthreeRunDeliberateStops{} deliberately, by the operator's interrupt of the first and by the wrapper's planned stop of the fourth, and \sthreeRunFailureStops{} by failures of the environment. Each later invocation resumed from the response cache with the same product commit, product settings, models, data and analysis files (\texttt{runs/\allowbreak{}study3/\allowbreak{}confirm/\allowbreak{}MANIFEST.json}). \texttt{RUNLOG.jsonl} records them as follows.
\begin{itemize}[leftmargin=1.2em,itemsep=1pt]
\item \emph{First} (\sthreeRunStartOne--\sthreeRunEndOne; stopped deliberately, not as planned). The driver logged ``interrupted during the checks'' and a spend of \sthreeRunFirstLoggedSpend. This line is wrong. A correction written by hand and dated \sthreeRunCorrectionDate{} \sthreeRunCorrectionTime{} states that all \sthreeRunChecks{} preflight checks had passed, that the invocation had completed Rules in context on both regimes and stored \sthreeRunFirstStoredOrig{} of \sthreeRunFirstExpectedOrig{} outputs of Rules in context with the original rules, spending USD~\sthreeRunFirstSpentOrig{} on that design, and that the operator interrupted it to resume with more parallel workers because the run was slow, not faulty. Two interrupt signals were sent; the second reached the driver while it was writing the log, so the wrong reason and spend were written. Spend is in the ledgers.
\item \emph{Second} (\sthreeRunStartTwo--\sthreeRunEndTwo; a failure), written by hand because the driver could not log it. It was killed by the time limit of the session's background command while Self-consistency was running on theft. The wrapper's interrupt at \sthreeRunWrapperMinutes{} minutes had no effect, because the driver had been started with the interrupt signal ignored.
\item \emph{Third} (\sthreeRunStartThree--\sthreeRunEndThree; a failure), written by hand for the same reason. The container that ran the session was restarted while \casepath was running on tenancy, with \sthreeRunThirdStoredTenancy{} of \sthreeRunThirdExpectedTenancy{} outputs stored; every process was killed without warning. The disk, the response cache and the product worktree were unchanged.
\item \emph{Fourth} (\sthreeRunStartFour--\sthreeRunEndFour, \sthreeRunFourthMinutes{} minutes; stopped deliberately, as planned). The driver logged ``interrupted'' while \casepath was running on theft, with \sthreeRunFourthStoredTheft{} of \sthreeRunFourthExpectedTheft{} outputs stored. A note written by hand and dated \sthreeRunCauseFourDate{} \sthreeRunCauseFourTime{} gives the cause: the operator's wrapper script sent one interrupt to the driver after \sthreeRunWrapperPlannedMinutes{} minutes so that the run would end before the session's time limit on background commands. It is the only planned early stop.
\item \emph{Fifth} (\sthreeRunStartFive--\sthreeRunEndFive). Complete; \casepath, checklist only, on the frontier model ran last, as planned.
\end{itemize}
Outputs are written after every case, so every stored output is complete. The first invocation used \sthreeRunWorkersFirst{} parallel workers and the later ones \sthreeRunWorkersLater. The main spending cap was lowered at each resumption to what remained of the planned USD~\sthreeRunCapOne{} (\sthreeRunCapTwo, \sthreeRunCapThree, \sthreeRunCapFour{} and \sthreeRunCapFive); the cap of USD~\sthreeRunCapFrontier{} for \casepath, checklist only, on the frontier model was unchanged. The ledgers record USD~\sthreeSpendTotal{} for the run: \sthreeSpendLuna{} on the cost-efficient model and \sthreeSpendTerra{} on the frontier model; \sthreeSpendFrontierProduct{} of it is that design and \sthreeSpendMain{} falls under the main cap.

\paragraph{Model access and settings}
Every design called its model through OpenRouter's chat-completions API. The product sent the identifiers \texttt{gpt-6-luna} and \texttt{gpt-6.1-terra} (\texttt{CASEPATH\_\allowbreak{}MODEL}); the harness sent them with the prefix \texttt{openai/} and asked OpenRouter for no fallback to another provider, for providers that accept every parameter and for providers that do not collect the data. The responses name \texttt{openai/\allowbreak{}gpt-6-luna} and \texttt{openai/\allowbreak{}gpt-6.1-terra}, served by OpenAI. No snapshot or version date is recorded. Every comparator call asked for low reasoning effort, with the reasoning text excluded from the response. The product ran with the commit's defaults plus \texttt{CASEPATH\_\allowbreak{}MODEL\_\allowbreak{}REASONING=true}: \sthreeReadings{} readings, the open rule \texttt{any} and full run scope (checklist scope for the checklist-only designs), with reasoning effort low for routing, the readings and the second document judgment and medium for the agents and the source checker. The other model-client settings were the commit's defaults: an output limit of 16,000 tokens and a timeout of 120 seconds. No design sent a temperature, top-p or seed; the provider's defaults applied and are not recorded. The manifest records the product's resolved settings for each product design and regime. The harness keys a cached response by the request's address, its full body and the run index. The three readings of a run differ in their response-schema name, so each is its own call; the full product and the checklist-only setting share a response whenever they send the same request in the same run.

\paragraph{Failed outputs}
\sthreeRunFailedOutputs{} outputs failed and were scored as empty checklists, as planned: \sthreeRunFailedReact{} from ReAct (\sthreeReactTenFailed{} tenancy, \sthreeReactThFailed{} theft), \sthreeRunFailedMer{} from Model executes rules, \sthreeRunFailedCore{} from the Core and \sthreeRunFailedSelfCons{} of the samples of Self-consistency (\sthreeRunFailedSelfConsSamplesTen{} tenancy, \sthreeRunFailedSelfConsSamplesTh{} theft). None came from a product design. A failed Self-consistency sample is an empty vote in the majority, not an empty checklist for the run. A post-hoc analysis leaves ReAct's failed outputs out (Appendix~\ref{app:rob-study3}).

\subsection{Before the run}
\begin{itemize}[leftmargin=1.2em,itemsep=1pt]
\item The run was started before the product's owners confirmed the merge of the tested commit, and at the time of writing they have not confirmed it. If they reject a change of behaviour, the test must be repeated on new sealed data.
%% AUTHORS: update the merge status of commit 6a72159c here and in Section 6.1.
\item Who wrote the texts of \texttt{fresh2/} and what that session could read is stated in Appendix~\ref{app:genai}. The session that directed development, drafted the plan and chose the open rule had written the test set's design and read its description and counts (\texttt{fresh2/\allowbreak{}README.md}), not its texts; before the freeze the texts were read only by the session that wrote them and by the blind-check sessions.
\item The probe claims were written with knowledge of the test set's design, its families and subsets, but not of its texts. The plan's words ``a writer that had not seen \texttt{fresh2/}'' hold for the texts only.
\item The order of the hypotheses was changed at 21:51 on 10 November, 38 minutes before the freeze (commit 482873e3). The first draft of the plan (commit 71eaa9eb, 19:17) had four hypotheses: theft against Rules in context on the same model; theft against every prior-work strategy; tenancy against the frontier model, as non-inferiority with margin 0.10; theft non-inferiority to the frontier model. The final draft (commit 984d2991, 21:26) had five: theft and tenancy against Rules in context on the same model; theft against every prior-work strategy; theft non-inferiority to the frontier model; tenancy against the frontier model. The reorder moved the last to third place. The plan's text at that commit says that the contrasts development supported most strongly come first. The work log at that commit records these development results: on theft stories 4--5, \casepath, checklist only, 0.40 wrong + missed changes per pair over nine runs, Rules in context 1.33, Self-consistency 1.50, ReAct 1.11 (one run), Model executes rules 2.06 and the frontier single prompt 0.67; on the 16 probe claims, 2.2 unsafe errors per denied-event claim for the product against 6.0--6.6 for the single prompts. The commit also adds development runs of the frontier single prompt on theft, on development data only. No design or scorer had touched \texttt{fresh2/}. Under the order of 21:26, the prior-work hypothesis would have been tested third and failed, the tenancy contrast against the frontier model would not have been tested, and only H1 and H2 would have been confirmed.
\item The default open rule (a doubt quoted by any reading blocks a presumption) was set after the probe and before the freeze, at 21:53 on 10 November, before any design read a sealed datum. An independent AI review session of the product's code found that with three readings a single ``not stated'' let a presumption override a doubt two readings had quoted. The directing author chose the rule \texttt{any} and thereby overruled a numerical criterion on wrong holds that the author had set earlier for this choice: on development data the rule caused about 0.09 more wrong holds per claim, and no premature request under either rule.
\item The comparators' texts, from the history of the harness files that hold them. Rules in context: the case-reading policy and task of Study~1 and the theft policy and task (\texttt{harness/\allowbreak{}tenancy\_\allowbreak{}methods.py}, \texttt{theft\_methods.py}) were first committed at 17:07 on 9 November, before Study~1's held-out runs, and never changed; the additions of Study~2 (\texttt{harness/\allowbreak{}v2/\allowbreak{}tenancy\_\allowbreak{}v2.py}, \texttt{theft\_v2.py}) were written at 22:22 that evening, and two of them changed at 23:12 during Study~2's development on development data, before Study~2's freeze at 00:15: the tenancy task's sentence on the silence policy, and the theft note on rules not admitted, which is left out when no rule is withheld. Model executes rules and ReAct for tenancy (\texttt{harness/\allowbreak{}product/\allowbreak{}baselines.py}) were first committed at 12:55 on 10 November and those for theft (\texttt{baselines\_theft.py}) at 19:10, during the product's development; each was run on development data before the freeze (\texttt{fresh/}) and neither file changed afterwards. Self-consistency sends Rules in context's messages five times; its sampling was committed at 19:10. Whether a text was edited before its first commit is not recorded. \casepath was revised over three studies on development data.
\item The driver's preflight checks passed at every invocation: the plan was committed unmodified; the product worktree was clean at the tested commit, with the expected source-tree hash and settings; the tenancy templates were identical to those of \texttt{fresh/}; and the product's evaluator agreed with the reference engines' exact deferral on every true, false and unresolved assignment of each tenancy workflow and on 3,000 random theft assignments (\texttt{harness/\allowbreak{}product/\allowbreak{}compare\_theft.py}; Appendix~\ref{app:proofs} reports a different check, of the product's evaluation against Equation~\eqref{eq:exact}, on 5,000 assignments).
\end{itemize}

\subsection{Departures from the plan's text}
We checked the plan sentence by sentence against what the driver and the analysis did. Every difference we found is listed.
\begin{itemize}[leftmargin=1.2em,itemsep=1pt]
\item \emph{Calls per case.} The plan's design table gives \sthreePlanCallsCasePath{} calls for \casepath and \sthreePlanCallsReact{} for ReAct. Measured: \sthreeCasePathTenCalls{} and \sthreeCasePathThCalls{} for \casepath, \sthreeReactTenCalls{} and \sthreeReactThCalls{} for ReAct (tenancy and theft). The checklist-only designs fall in the stated ranges.
\item \emph{Cost of Self-consistency.} The plan states that Self-consistency costs at least as much as \casepath, checklist only, on both regimes, and treats the checklist-only row as the comparison matched in cost. This holds on tenancy (USD~\sthreeSelfConsTenCostFour{} against \sthreeChecklistTenCostFour) but not on theft (\sthreeSelfConsThCostFour{} against \sthreeChecklistThCostFour), so no design is matched in cost.
\item \emph{Full product and checklist only share responses.} Identical calls are served from the response cache. \sthreeSharedTenCalls{} of \sthreeSharedTenAllCalls{} tenancy calls and \sthreeSharedThCalls{} of \sthreeSharedThAllCalls{} theft calls of the full product (USD~\sthreeSharedTenUSD{} and \sthreeSharedThUSD) returned responses first paid by the checklist-only run, among them routing and some readings. The contrast between the two is therefore not between independent runs. The plan does not address this.
\item \emph{Core runs share readings.} The Core's driver pairs each run's reading with the next run's reading, cycling through the runs, so its runs are not independent; its cost per text-run is \sthreeCoreMultiplier{} times the cost of the readings paid.
\item \emph{How cost was computed.} The plan says cost per case from the ledger. For the product designs, each distinct call of a stored case-run was joined to the ledger by its cache key, whoever paid it, so the shared calls above count at full price. Other designs: matching ledger rows of the cases in scope, each call once. The split by agent needs the stored agent list to line up with the calls; it did not in \sthreeCasePathTenNotLinedUp{} of \sthreeCasePathTenCaseRuns{} tenancy and \sthreeCasePathThNotLinedUp{} of \sthreeCasePathThCaseRuns{} theft case-runs, whose cost is reported as not attributed. The product books the brief and the first reading to the claim reader and the further readings and second document judgments to the source checker, so the registered split counts the brief among the agents that can decide the checklist; Appendix~\ref{app:rob-study3} reconciles it with the cost of the checklist-only setting.
\item \emph{Subsets.} The plan defines the subsets by the defining event; the analysis selects them from the reference: denied event, an empty request set and no held document; open event, an empty request set and a held document; ordinary, a non-empty request set. The counts equal the plan's (\sthreeClaimsDenied, \sthreeClaimsOpenEvent, \sthreeClaimsOrdinary; \sthreeClaimsOpenBranch{} with an open branch fact), and the analysis checks that the three form a partition.
\item \emph{Subset rows beside H2 and H3.} The plan says the subsets are reported beside each tenancy hypothesis. The analysis reports the ordinary claims and the denied- and open-event claims together, as contrasts over the families with at least one claim in the subset (\sthreeHTwoOrdinaryClusters{} and \sthreeHTwoDeniedOpenClusters{} families), with interval and $p$ marked descriptive. Their values are means of family values; per-claim values are in Table~\ref{tab:s3-subsets}.
\item \emph{Number of tests.} The procedure has \sthreeHypotheses{} hypotheses and \sthreePrimaryTests{} tests (H4 has \sthreeHFourTests). Beside them the analysis reports \sthreeBesideContrasts{} descriptive subset contrasts and \sthreeSecondaryContrasts{} secondary contrasts. Of these, \sthreeSecondaryVsCasePath{} (\casepath on the frontier model against \casepath, both regimes) are not listed in the plan; they are descriptive.
\item \emph{H5.} Not tested. The plan asks for its non-inferiority $p$; it is in Table~\ref{tab:s3-intervals}, labelled outside the procedure, with a caveat on the bootstrap $p$ at nine clusters.
\item \emph{Resampling.} The plan fixes \sthreeBootstrap{} resamples but not the seed; the specification fixes it (\sthreeSeed).
\item \emph{Agreement and repeated errors.} Run-to-run agreement is the Jaccard index of the requested documents over pairs of runs; holds are not included, and the one-run design has none. Repeated errors counted once are reported as values, without tests.
\item \emph{Two-valued on theft.} With unresolved taken as false, the Two-valued design gives the same theft requests as \casepath, and theft scores requests only, so the two rows are equal.
\item \emph{Settings on resumption.} The plan says an interrupted run resumes with settings unchanged. Product, model and analysis settings were unchanged; the number of workers and the main cap (see above) were changed.
\end{itemize}

\subsection{Analyses added after review}
\label{app:review-analyses}
After an independent referee report and a claims audit we added eleven analyses, and after a second referee report two more (\texttt{analysis/\allowbreak{}study3\_\allowbreak{}review.py}, output \texttt{analysis/\allowbreak{}study3\_\allowbreak{}review.md}). None is pre-registered, none is part of the plan, and none changes a hypothesis's status. Each answers a question of the review.
\begin{itemize}[leftmargin=1.2em,itemsep=1pt]
\item \emph{R1, claim-level outcomes.} The per-document measure can count one misjudged gating fact as several errors; R1 counts whole checklists.
\item \emph{R2, other mixes of claim types.} The denied- and open-event claims were added by design; R2 reweights the subsets to other prevalences.
\item \emph{R3, fact-level accuracy.} Proposition~\ref{prop:errors} needed a measured rate of wrong admitted values, including wrong values on which all readings agreed.
\item \emph{R4, wrong presumptions.} The texts contain no unmentioned fact that is true; R4 makes some true and rescores the stored checklists.
\item \emph{R5, theft inference.} With nine clusters the percentile bootstrap under-covers and stories are crossed with facts; R5 gives inverted sign-flip intervals, stories as clusters and two-way resampling.
\item \emph{R6, family-level differences.} The tenancy tests rest on 24 families, nine of which hold the denied- and open-event claims; R6 shows each family.
\item \emph{R7, full product against checklist only.} The comparison shares cached responses and its cost split was inconsistent; R7 counts identical checklists and reconciles the cost.
\item \emph{R8, single prompts on denied-event claims.} R8 asks whether the single prompts misread the denial or requested documents despite it.
\item \emph{R9, reviewed against original rules.} The reviewed rules raised the single prompts' errors; R9 looks for the cause.
\item \emph{R10, time per case.} Latency was not reported.
\item \emph{R11, exact deferral against two-valued and strong Kleene evaluation.} R11 checks where the third value changed a checklist.
\item \emph{R7b, run-to-run baseline} (added after the second review round). R7b counts how often two runs of the same setting gave identical checklists, so that R7's shares can be read against it.
\item \emph{R11b, strong Kleene evaluation scored} (added after the second review round). R11 showed that exact deferral changed checklists; R11b scores both evaluations against the reference.
\end{itemize}
